% v9, corrected against First Proof paper (Cappell-Weinberger-Yan, "Fowler's theorem for involutions")

\documentclass[11pt]{article}

\usepackage[margin=1in]{geometry}
\usepackage{amsmath, amssymb, amsthm}
\usepackage{mathtools}
\usepackage{xcolor}
\usepackage{hyperref}
\hypersetup{colorlinks=true,
            linkcolor=blue!50!black,
            citecolor=blue!50!black,
            urlcolor=blue!50!black}

%---- math macros ----
\newcommand{\R}{\mathbb{R}}
\newcommand{\Q}{\mathbb{Q}}
\newcommand{\Z}{\mathbb{Z}}
\newcommand{\Lgp}{\mathbb{L}}
\newcommand{\sign}{\operatorname{sign}}
\newcommand{\cd}{\operatorname{cd}}
\newcommand{\tM}{\widetilde M}
\DeclareMathOperator{\im}{im}

%---- theorem environments ----
\newtheorem{theorem}{Theorem}
\newtheorem{lemma}{Lemma}
\newtheorem{proposition}{Proposition}

\title{Uniform Lattices with $2$-Torsion as Fundamental Groups}
\author{}
\date{May 27, 2026}

\begin{document}
\maketitle

\section*{Statement and Answer}

Let $\Gamma$ be a uniform (cocompact) lattice in a real semisimple Lie group
$G$, and suppose $\Gamma$ contains an element of order $2$.

\medskip
\noindent\textbf{Theorem.} \emph{Such a $\Gamma$ cannot be the fundamental group
of a closed manifold without boundary whose universal cover is acyclic over
$\Q$.}

\medskip
\noindent The answer is \textbf{NO}.

\medskip
A torsion-free uniform lattice in $G$ \emph{is} the fundamental group of such a
manifold --- a closed locally symmetric space, whose universal cover is the
contractible symmetric space $X=G/K$. The content of the theorem is that
$2$-torsion obstructs this. Fowler, in his Ph.D.\ thesis, proved the
corresponding statement when $\Gamma$ has \emph{odd} torsion (a proof appears in
\cite{W2}); the present argument extends Fowler's theorem to the case of
$2$-torsion, following Cappell--Weinberger--Yan.

We stress at the outset what does \emph{not} work: a proof using only that $M$
is a finite Poincar\'e complex over $\Q$ cannot succeed, because such complexes
with these fundamental groups exist (Proposition~\ref{prop:counterexample}). The
argument must use that $M$ is a genuine manifold, and it does so through the
symmetric signature and an equivariant fundamental-class computation.

\section*{What a finite-complex / Poincar\'e-duality argument cannot see}

It is tempting to argue purely with rational Poincar\'e duality and Euler
characteristics. That route is blocked by an explicit construction of Fowler.

\begin{proposition}\label{prop:counterexample}
There exist groups of the form $\pi_1(M_3)\times\Gamma$, with $\Gamma$ a lattice
with torsion, that are fundamental groups of spaces having the rational homotopy
type of a finite complex satisfying rational Poincar\'e duality.
\end{proposition}

\begin{proof}[Construction]
Let $\Gamma$ be a lattice in a linear semisimple group $G$, and choose a
homomorphism $\Gamma\to\Delta$ to a finite group with torsion-free kernel
$\Gamma_0$. Let $E\Delta$ be a contractible free $\Delta$-space and $M_3$ any
closed hyperbolic $3$-manifold. The space
\[
  M_3\times\bigl(K\backslash G/\Gamma_0\times E\Delta\bigr)/\Delta
\]
has the rational homotopy type of a finite complex and satisfies rational
Poincar\'e duality; its fundamental group is $\pi_1(M_3)\times\Gamma$, a lattice
in $\mathrm{SO}(3,1)\times G$.
\end{proof}

In particular, ``multiplicativity of the Euler characteristic in finite covers''
\emph{fails} for infinite complexes with finitely generated rational homology: a
rationally acyclic example with rational Euler characteristic $1$ is the
universal cover of $\R P^2\vee(\bigvee_\infty S^2)$ after equivariantly attaching
cells $D^3\times\Z/2$ along a free $\Z[1/2][\Z/2]$-basis of
$\pi_2=\Z[-1]\oplus\Z[\Z/2]^{\infty}$ to kill the homology; both it and its
$\Z/2$-quotient become rationally acyclic with rational Euler characteristic $1$.
Hence the obstruction cannot be an Euler-characteristic denominator, and any
proof relying on $\chi(M)\in\Z$ versus $\chi_\Q(\Gamma)\notin\Z$ is invalid. We
also note that any proof resting on a naive Lefschetz fixed-point argument is
unsound: the relevant Lefschetz-type statement is \emph{false} (witness
$\R^1$ with $f$ a translation, which has no fixed points yet Lefschetz number
$-1$ in that sense).

\section*{Proof of the Theorem}

\subsection*{Reduction}

Passing to a normal subgroup of finite index, it suffices to treat
$\Gamma=\pi\rtimes\Z_2$ for a torsion-free $\pi$, a lattice in $G$, such that
there is an involution on $M_0:=K\backslash G/\pi$ --- by isometries of the
locally symmetric metric --- whose fixed set $F$ is nonempty. (If $F$ is
disconnected one argues componentwise; we write the connected case.) Here $M_0$
is a closed aspherical manifold, a $K(\pi,1)$.

Suppose, for contradiction, that $X^m$ is a closed manifold with
$\pi_1(X)\cong\Gamma$ and rationally acyclic universal cover. Let $Y\to X$ be the
$2$-fold cover corresponding to $\pi\le\Gamma$; then $\pi_1(Y)\cong\pi$ and the
universal cover of $Y$ (the same as that of $X$) is rationally acyclic.

\subsection*{Step 1: a degree-one fundamental class}

Work with symmetric signatures in the symmetric $L$-group $L(\R\pi)$ over the
reals. (Over $\R$ the symmetric and quadratic $L$-groups agree, so we may write
$L(\R\pi)$ unambiguously.) There is a map $f:Y\to M_0$ inducing the identity on
$\pi_1=\pi$. It need not have degree one, but over $\R$ every degree has a square
root, so the symmetric signature is sensitive only to the sign of the degree,
and $f$ induces an equivalence of symmetric signatures.

Because the Novikov conjecture holds for $\pi$ (a lattice in a semisimple group;
it satisfies the Borel/Novikov package via its proper action on the
nonpositively curved $X$), the assembly map
\[
  H_m\bigl(B\pi;\Lgp(\R)\bigr)\;\longrightarrow\;L_m(\R\pi)
\]
is injective. The symmetric signature of $Y$ is detected in the degree-$m$ piece
$H_m(B\pi;\Z)$ as the image of the fundamental class. Injectivity forces
\begin{equation}\label{eq:deg-one}
  f_*[Y]\;=\;[M_0]\quad\text{in}\quad H_m(B\pi;\Q),
\end{equation}
i.e.\ $f$ is rationally degree one onto the aspherical $M_0$.

\subsection*{Step 2: an equivariant cobordism and the fixed set}

We now run the cobordism argument of \cite{W1}. Consider any closed manifold $Z$
with fundamental group $\pi$ carrying an involution inducing the given
automorphism of $\pi$, together with the image of $[Z]$ in $H_m(B\Gamma;\Z_2)$.
Standard equivariant homotopy theory gives an equivariant map $g:Z\to M_0$,
hence a map of fixed sets $Z^{\Z_2}\to F$. We claim
\begin{equation}\label{eq:fixed}
  g_*[Z]\;=\;g_*[Z^{\Z_2}],
\end{equation}
using the map $\Z_2\times\pi_1F\to\Gamma$ and the periodicity in the group
homology of $\Z_2$ to raise the dimension from $\dim F$ to $\dim M_0$.

The relevant cobordism realizing \eqref{eq:fixed} runs between $Z$ and a
projective-space bundle over $Z^{\Z_2}$ --- the projectivized normal bundle of
$Z^{\Z_2}$ --- whose fundamental class is the asserted element by the
Leray--Hirsch theorem. Explicitly it is $Z\times[0,1]$, with $Z\times\{1\}$
modified by quotienting the free $\Z/2$-action on the complement of an
equivariant regular neighbourhood of $Z^{\Z_2}$.

\subsection*{Step 3: the contradiction}

Apply Step 2 in two ways.

\emph{For $Y$.} The deck $\Z_2$-action on $Y$ (equivalently the action with
quotient $X$) is \emph{free}, so $Y^{\Z_2}=\varnothing$ and the right-hand side
of \eqref{eq:fixed} is $0$. Hence the corresponding image of the fundamental
class is $0$.

\emph{For $M_0$.} The isometric $\Z_2$-action on $M_0$ has fixed set $F$ which is
aspherical, and the Borel construction shows that
$\Z_2\times F\to\Gamma$ induces an injection on homology in dimension
$\dim(M_0/\Z_2)$ (an isomorphism in higher dimensions; see \cite{Borel}). Since
the fundamental class of an aspherical manifold is nontrivial in its group
homology, the corresponding image for $M_0$ is \emph{nonzero}.

But \eqref{eq:deg-one} identifies the $Y$-image with the $M_0$-image rationally:
the degree-one map $f$ carries the (zero) class for $Y$ to the (nonzero) class
for $M_0$. This is a contradiction. Therefore no such closed manifold $X$
exists, and $\Gamma$ is not the fundamental group of a closed manifold with
rationally acyclic universal cover.\qed

\section*{Remarks}

\begin{enumerate}\setlength\itemsep{4pt}
\item \textbf{Where the manifold hypothesis enters.} The argument uses that $X$
is a manifold twice: to form the symmetric signature and apply assembly
injectivity (Step 1), and to run the equivariant normal-bundle cobordism (Step
2). Proposition~\ref{prop:counterexample} shows neither can be replaced by
``finite Poincar\'e complex'' alone; any proof using only finite complexes and
Poincar\'e duality must fail.
\item \textbf{Odd torsion.} The corresponding statement for odd torsion is
Fowler's theorem \cite{W2}; the present argument adapts the involution case
($p=2$), where the fixed-point and normal-bundle analysis replaces the
Euler-characteristic counting available in nicer settings.
\item \textbf{Role of Novikov.} Only injectivity of the rational assembly map
for $\pi$ is needed, which holds for lattices in semisimple groups via their
proper actions on the nonpositively curved symmetric space.
\item \textbf{Honest accounting.} The delicate points are
\eqref{eq:deg-one}\,--\,\eqref{eq:fixed}: identifying the symmetric signature with
the fundamental class under assembly, and the periodicity/Leray--Hirsch
computation that pushes the fixed-set class up to top dimension. These are the
substance of the surgery-theoretic argument and are stated here at the level of
their structure.
\end{enumerate}

\begin{thebibliography}{9}
\bibitem{Borel} A.~Borel.
  \emph{Seminar on Transformation Groups}. Princeton University Press, 1960.
\bibitem{W1} S.~Weinberger.
  Group actions and higher signatures II.
  \emph{Comm. Pure Appl. Math.}, 1987.
\bibitem{W2} S.~Weinberger.
  \emph{Variations on a Theorem of Borel}.
  Cambridge University Press, 2022.
\end{thebibliography}

\end{document}